\subsection{星代数中的正元素}

定义 6. --- 设 A 是一个星代数。称 A 的元素 x 是正的，如果它是 Hermite 的且 \(\mathrm{Sp}_{\mathrm{A}}^{\prime}(x) \subset \mathbf{R}_{+}\)。用 \(\mathrm{A}_{+}\) 表示 A 的正元素组成的集合。它是 \(\mathrm{A}_{h}\) 的一个子集。

我们记 \(x \geqslant y\)，如果 \(x - y \in \mathrm{A}_{+}\)。

若星代数 A 是幺的，则它的单位元是正的。

若 B 是 A 的一个星子代数，则有 \(B_{+} = B \cap A_{+}\) (I, p.~106 命题 5 的推论)。

若 \(\pi\colon A\to B\) 是星代数的一个态射，则 \(\pi(A_{+})\subset B_{+}\)。

例. --- 1) 设 X 是一个局部紧空间。在 X 上在无穷远处趋于 0 的连续函数组成的星代数 \(\mathcal{C}_{0}(\mathrm{X})\) 中，相应地在 X 上有界连续函数组成的星代数 \(\mathcal{C}_{b}(\mathrm{X})\) 中，函数 \(f\) 是正元素当且仅当它取实值且 \(f(x) \geqslant 0\) 对一切 \(x \in \mathrm{X}\) 成立（参见 I, p.~17, 例 3）。

\begin{enumerate}
\def\labelenumi{\arabic{enumi})}
\setcounter{enumi}{1}
\item
  设 A 是一个交换星代数。设 \(a\) 是 A 的一个元素。由于 \(\mathrm{Sp}_{\mathrm{A}}^{\prime}(x)\) 是 \(\{0\}\) 与 Gelfand 变换 \(\mathcal{G}(a)\) 的像的并 (I, p.~37, 命题 6)，元素 \(a\) 是正的当且仅当 Gelfand 变换 \(\mathcal{G}(a)\) 是一个正函数。
\item
  设 E 是一个复 Hilbert 空间。星代数 \(\mathcal{L}(\mathrm{E})\) (I, p.~102, 例 1) 的元素 x 是正的当且仅当它是 EVT, V, p.~45, 定义 6 意义下 E 的正自同态 (I, p.~138, 命题 8)。
\end{enumerate}

引理 12. --- 设 A 是一个幺星代数，\(x \in A\) 是一个 Hermite 元。

\begin{enumerate}
\def\labelenumi{\alph{enumi})}
\item
  元素 \(x\) 是正的当且仅当 \(\left\| \|x\| \cdot 1 - x \right\| \leqslant \|x\|\)；
\item
  若 \(\| x\| \leqslant 1\)，则 \(x\) 是正的当且仅当 \(\| 1 - x\| \leqslant 1\)。
\item
  若 x 是正的，则 1 - x 是正的当且仅当 \(\|x\| \leqslant 1\)；
\item
  若 x 是正的且 \(y \in A_{+}\) 与 x 交换，则 xy 是正的。
\end{enumerate}

元素 \(x\) 是 Hermite 的，因而是正规的。通过考虑由 \(x\) 生成的交换星子代数，可化归为代数 A 是交换的情形，即对某个局部紧拓扑空间 X 有 \(A = \mathcal{C}_0(X)\) 的情形 (I, p.~108, 定理 1)。于是前三个断言由上面的例 1 立即得出。同样，为证断言 \(d\))，可以考虑由 \(x\) 与 \(y\) 生成的星子代数，它是交换的。

命题 14. --- 设 A 是一个星代数。集合 \(\mathrm{A}_{+}\) 是实 Banach 空间 \(\mathrm{A}_{h}\) 中的一个闭的尖凸锥 (EVT, II, p.~11)。

设 \(\widetilde{\mathsf{A}}\) 是由 \(\mathsf{A}\) 添加单位元而得的星代数。由于 \(\mathsf{A}_{+} = \mathsf{A} \cap \widetilde{\mathsf{A}}_{+}\)，只需对 \(\widetilde{\mathsf{A}}\) 证明该命题。于是可设 \(\mathsf{A}\) 有单位元。

我们有 \(0 \in \mathrm{A}_{+}\)。对每个 \(\lambda \in \mathbf{R}_{+}^{*}\) 及每个 \(x \in \mathrm{A}\)，有 \(\mathrm{Sp}_{\mathrm{A}}^{\prime}(\lambda x) = \lambda \mathrm{Sp}_{\mathrm{A}}^{\prime}(x)\)，这蕴含 \(\mathrm{A}_{+}\) 是实 Banach 空间 \(\mathrm{A}_{h}\) 中的一个锥。

为证 \(A_{+}\) 是凸的，只需证明若 x 与 y 是正的，则 \(x + y \geqslant 0\) (EVT, II, p.~11, 命题 10)。通过位似，只需证明：若 \(x \geqslant 0\) 与 \(y \geqslant 0\) 还满足 \(\|x\| \leqslant 1\)、\(\|y\| \leqslant 1\)，则元素 \(\frac{1}{2}(x + y)\) 是正的。但我们有

\[
\left\| 1 - \frac {1}{2} (x + y) \right\| \leqslant \frac {1}{2} \| 1 - x \| + \frac {1}{2} \| 1 - y \| \leqslant 1
\]

依引理 12 的断言 \(b)\)，而这同一断言表明 \(\frac{1}{2}(x + y)\) 是正的。

最后，引理 12 的断言 \(a\)) 也蕴含 \(\mathrm{A}_{+}\) 是闭的。

由于 \(A_{+}\) 是 \(A_{h}\) 中的一个以 0 为顶点的锥，它是突出的当且仅当 \(\mathrm{A}_{+}\cap(-\mathrm{A}_{+})\) 缩减为 0。但若 \(x\in\mathrm{A}_{+}\cap(-\mathrm{A}_{+})\)，则有 \(\mathrm{Sp}_{\mathrm{A}}^{\prime}(x)=\{0\}\)，于是 \(\varrho(x)=0\)，且 \(\|x\|=0\)。由于 x 是 Hermite 的 (I, p.~104, 引理 7, (2))，由此得 x=0。

命题 14 表明关系 \(x\geqslant y\) 是 \(\mathrm{A}_h\) 上的一个序关系 (EVT, II, p.~13, 命题 13)。

命题 15. — 设 A 是一个星代数。设 x 是 A 的一个正规元。

\begin{enumerate}
\def\labelenumi{\alph{enumi})}
\item
  设 A 是幺的，f 是从 \(\mathrm{Sp}_{\mathrm{A}}(x)\) 到 C 中的一个连续函数。为使 \(f(x)\) 是正的，必须且只需 f 的像包含在 \(R_{+}\) 中；
\item
  设 f 是从 \(\mathrm{Sp}_{\mathrm{A}}^{\prime}(x)\) 到 C 中的一个连续函数，满足 \(f(0)=0\)。为使 \(f(x)\) 是 A 的一个正元素，必须且只需 f 的像包含在 \(R_{+}\) 中。
\end{enumerate}

断言 \(a\)) 由 I, p.~111 推论 2 的断言 \(a\)) 得出，断言 \(b\)) 由 I, p.~114 的命题 13 得出。

设 \(x\) 是星代数 A 的一个 Hermite 元。它的谱包含在 \(\mathbf{R}\) 中 (I, p.~106, 命题 4)。考虑从 \(\mathrm{Sp}_{\mathrm{A}}^{\prime}(x)\) 到 \(\mathbf{R}\) 中的由下式定义的连续函数

\[
f _ {1} \colon t \mapsto \sup (t, 0), \quad f _ {2} \colon t \mapsto \sup (- t, 0), \quad f _ {3} \colon t \mapsto | t |.
\]

我们记

\[
x ^ {+} = f _ {1} (x), \quad x ^ {-} = f _ {2} (x), \quad | x | = f _ {3} (x).\tag{3}
\]

由于函数 \(f_1, f_2, f_3\) 取实值、是正的且在 0 处为零，元素 \(x^+, x^-\) 与 \(|x|\) 是 A 的正元素（命题 15, a)），且属于 A 的由 \(x\) 生成的星子代数。

我们有 \(f_{1}(t) - f_{2}(t) = t\) 对一切 \(t \in \mathbf{R}\) 成立，且有关系式 \(f_{1} + f_{2} = f_{3}\) 与 \(f_{1}f_{2} = 0\)。由此得

\[
x = x ^ {+} - x ^ {-}, \quad | x | = x ^ {+} + x ^ {-}, \quad x ^ {+} x ^ {-} = x ^ {-} x ^ {+} = 0.\tag{4}
\]

由于函数演算映射是等距的，我们有

\[
\| | x | \| = \| x \|, \quad \| x ^ {+} \| \leqslant \| x \|, \quad \| x ^ {-} \| \leqslant \| x \|.
\]

设 \(x\) 是 A 的一个正元素。它是 Hermite 的，因而是正规的。设 \(\alpha \in \mathbf{R}_{+}^{*}\)，并设 \(g\) 是 \(\mathrm{Sp}_{\mathrm{A}}^{\prime}(x)\) 上函数 \(t \mapsto t^{\alpha}\) 的限制；我们记 \(x^{\alpha} = g(x)\)。这是 A 的由 \(x\) 生成的星子代数中的一个正元素。设 \(\alpha\) 与 \(\beta\) 属于 \(\mathbf{R}_{+}^{*}\)。由于函数演算映射是一个代数同态，并依 I, p.~114 的命题 13，我们有

\[
x ^ {\alpha} x ^ {\beta} = x ^ {\alpha + \beta} \quad (x ^ {\alpha}) ^ {\beta} = x ^ {\alpha \beta}.\tag{5}
\]

我们也记作 \(\sqrt{x} = x^{1/2}\)。

命题 16. — 设 x 是 A 的一个正元素。设 \(\alpha \in \mathbf{R}_{+}^{*}\)。那么 \(x^{1/\alpha}\) 是 A 中满足 \(y^{\alpha} = x\) 的唯一的正元素 y。

我们在上面已看到 \(x^{1/\alpha}\) 满足所要求的性质。反之，设 y 是 A 中满足 \(y^{\alpha} = x\) 的一个正元素。依公式 (5)，我们有 \(y = (y^{\alpha})^{1/\alpha} = x^{1/\alpha}\)，这就是要证的。

引理 13. — 设 A 是一个幺星代数。A 的每个元素都是酉元之和。

设 \(x\) 是 A 的一个 Hermite 元。首先设 \(\| x \| \leqslant 2\)。依引理 12, c)，我们有 \(1 - \frac{1}{4} x^2 \in \mathrm{A}_+\)。令 \(y = \frac{1}{2} x + i\sqrt{1 - \frac{1}{4} x^2}\)。我们有 \(y^* = \frac{1}{2} x - i\sqrt{1 - \frac{1}{4} x^2}\)，于是 \(yy^* = 1\)，且 \(x = y + y^*\) 是两个酉元之和。在一般情形，取整数 \(k\) 使 \(\| \frac{1}{k} x \| \leqslant 2\)；于是元素 \(x\) 是 \(2k\) 个酉元之和。依 I, p.~96 的引理 2，引理得证。

定理 2. — 设 A 是一个星代数。A 的元素 x 是正的当且仅当存在 \(y \in A\) 使 \(x = y^{*}y\)。

设 \(x\) 是正的。令 \(y = x^{1/2}\)；它是 A 的一个 Hermite 元，且我们有 \(y^*y = y^2 = x\)。

反之，设 y 是 A 的一个元素，并令 \(x = y^{*}y\)。这是 A 的一个 Hermite 元。我们来证明 x 是正的。为此，记 \(z = x^{-}\) 并令 w = yz。于是我们有

\[
w ^ {*} w = z ^ {*} y ^ {*} y z = z x z = z (x ^ {+} - z) z = - z ^ {3}.
\]

由于 \(z \geqslant 0\)，有 \(z^3 \geqslant 0\)，于是我们推出 \(\mathrm{Sp}_{\mathrm{A}}'(w^{*}w) \subset \mathbf{R}_{-}\)。再写 \(w = w_1 + iw_2\)，其中 \(w_1\) 与 \(w_2\) 是 Hermite 的 (I, p.~96, 引理 2)。元素 \(w_1^2\) 与 \(w_2^2\) 是正的。我们有 \(ww^{*} + w^{*}w = 2w_{1}^{2} + 2w_{2}^{2}\)，于是命题 14 表明 \(ww^{*} = 2w_{1}^{2} + 2w_{2}^{2} + (-w^{*}w)\) 是正的。由于 \(\mathrm{Sp}_{\mathrm{A}}'(ww^{*}) = \mathrm{Sp}_{\mathrm{A}}'(w^{*}w)\) (I, p.~5, 命题 1)，而后者包含在 \(\mathbf{R}_{-}\) 中，我们得出结论 \(\mathrm{Sp}_{\mathrm{A}}'(w^{*}w) = \{0\}\)。由于 \(w^{*}w\) 是 Hermite 的，这蕴含 (I, p.~108, 推论 1) \(\|w^{*}w\| = \varrho(w^{*}w) = 0\)，从而 \(z^{3} = 0\)。由于 z 是 Hermite 的，我们有 z = 0。于是 \(x = x^{+}\) 是正的。

注记. --- 设 A 是一个星代数，\(x \in A\)。元素 \(x^{*}x\) 是正的；因此我们令 \(|x| = (x^{*}x)^{1/2}\)。我们有 \(\|x\|^{2} = \|x^{*}x\| = |||x|^{2}\| = |||x||^{2}\)，从而 \(\|x\| = |||x|\|\)。

当 x 是正规元时，我们还有 \(|x| = f(x)\)，其中 f 是从 C 到 C 的在 0 处为零、由 \(f(z) = |z|\) 给定的映射。特别地，当 x 是 Hermite 元时，\(|x|\) 与公式 (3) 定义的元素重合。

进一步设 A 是幺的且 x 是可逆的。那么 \(|x|\) 也是可逆的，元素 \(u = x|x|^{-1}\) 是酉元，且我们有 \(x = u|x|\)（极分解；关于 Hilbert 空间自同态的情形另见 I, p.~139, 第 8 小节）。

引理 14. --- 设 A 是一个星代数，x 与 y 是 A 的 Hermite 元。

\begin{enumerate}
\def\labelenumi{\alph{enumi})}
\item
  若 \(x \leqslant y\)，则对 A 的每个元素 \(w\)，有 \(w^{*}xw \leqslant w^{*}yw\)。特别地，若 \(y \geqslant 0\)，则有 \(w^{*}yw \geqslant 0\)；
\item
  设 A 是幺的。若 \(0 \leqslant x \leqslant y\) 且 \(x\) 可逆，则 \(y\) 可逆且 \(y^{-1} \leqslant x^{-1}\)；
\item
  若 \(0 \leqslant x \leqslant y\)，则 \(\|x\| \leqslant \|y\|\)。
\end{enumerate}

令 \(u = (y - x)^{1/2}\)。我们有 \(w^{*}yw - w^{*}xw = w^{*}u^{2}w = (uw)^{*}(uw)\)，于是断言 a) 由定理 2 得出。

我们来证明断言 \(b\))。首先设 \(x = 1\)。设 B 是由 \(y\) 生成的幺星子代数。依 Gelfand 同构，\(y\) 对应于一个 \(\geqslant 1\) 的连续函数，定义在紧空间 \(\mathsf{X}(\mathsf{B})\) 上。于是这个函数是可逆的，且它的逆 \(\leqslant 1\)。这蕴含 \(y\) 可逆且 \(y^{-1} \leqslant 1 = x^{-1}\)。在一般情形，我们注意到 \(0 \leqslant 1 \leqslant x^{-1/2}yx^{-1/2}\)（依 \(a\))），故前面的情形蕴含 \(z = x^{-1/2}yx^{-1/2}\) 可逆且 \(z^{-1} \leqslant 1\)。于是 \(y\) 可逆且 \(y^{-1} \leqslant x^{-1}\)，这仍依 \(a\))。

为证断言 \(c\))，可设 A 是幺的 (I, p.~105, 命题 3)。首先设 \(y\) 可逆。令 \(b = \sqrt{y}\)。依 \(a\))，条件 \(0 \leqslant x \leqslant y\) 蕴含 \(0 \leqslant b^{-1}xb^{-1} \leqslant b^{-1}yb^{-1} = 1\)，于是 \(\|b^{-1}xb^{-1}\| \leqslant 1\)，依 I, p.~116 的引理 12 \(c\))。于是我们有

\[
\| x \| = \| b (b ^ {- 1} x b ^ {- 1}) b \| \leqslant \| b \| \| b ^ {- 1} x b ^ {- 1} \| \| b \| \leqslant \| b \| ^ {2} = \| b ^ {2} \| = \| y \|.
\]

在一般情形，对每个实数 \(\varepsilon > 0\)，元素 \(y + \varepsilon\) 可逆且 \(0 \leqslant x \leqslant y + \varepsilon\)。于是由前所证，\(\|x\| \leqslant \|y + \varepsilon\|\) 对每个实数 \(\varepsilon > 0\) 成立，由此得结果。

注. — 一般地，若 x 与 y 是星代数 A 的正元素，条件 \(0 \leqslant x \leqslant y\) 并不蕴含 \(x^{2} \leqslant y^{2}\)（参见 I, p. 182, 习题 15）。

